\section{Avoiding Documentation Drift}

The foundations chapter identified why documentation drift is structural rather than accidental: code changes are continuously exercised by compilers, automated tests, and continuous integration pipelines, so problems introduced into the code tend to be caught quickly, while a specification file has no built-in awareness that the code it describes has changed, and is instead updated manually, opportunistically, and often incompletely. This framing matters for what kind of practice actually prevents drift, because it rules out relying on manual diligence alone --- the same absence of automatic checking that allows drift to occur in the first place is exactly what manual diligence has no structural way to compensate for consistently.

Part 1's OpenAPI specification had no automated check tying it to the implementation at all, which means any divergence between what the specification claimed and what the routes actually did would only be noticed if someone happened to compare the two directly --- the same failure mode the foundations chapter describes as the general cause of documentation drift, present here in its plainest form. \textbf{The practice this rules out is treating documentation accuracy as something to be maintained by discipline or convention alone; without an automatic check, a specification's correctness depends entirely on someone remembering to update it every time, which is precisely the condition under which drift accumulates unnoticed.}

Part 2 addressed this not by generating the specification automatically from the code, but by keeping the specification hand-written and adding a dedicated pytest module that inspects the Flask application's routing table at runtime and compares it against the specification file, asserting that every implemented route and method is documented and that no documented endpoint is missing from the implementation. \textbf{The recommendation this supports is to fold a consistency check between code and documentation into the same automated pipeline that already checks the code itself, rather than either generating documentation automatically or trusting it to manual updates --- a mismatch then produces a failing test, using the same enforcement mechanism already relied on for correctness elsewhere, instead of depending on someone to notice a discrepancy on their own.}

This choice also reflects a trade-off worth stating explicitly rather than treating the decision as automatic. Generating a specification directly from source-code annotations guarantees consistency by construction, but requires maintaining that metadata on every endpoint and ties the codebase more tightly to the generation tooling; keeping the specification as a separate, hand-written artifact avoids that coupling, at the cost of needing an explicit check to catch what automatic generation would have prevented from arising at all. \textbf{The practice this points to is not that one approach is unconditionally correct, but that whichever approach is chosen, the choice should be paired with a deliberate mechanism --- either generation-by-construction or an automated comparison --- rather than left as an assumption that a document and the code it describes will simply stay in agreement on their own.}

